% !TeX root = ../../Book2.tex
\chapter{常微分方程基础}
\label{b2:ch:ode}
\BookChapterNumber{b2:ch:ode}
\begin{chapterabstract}
常微分方程用当前位置和时刻决定瞬时变化率，再由初始数据确定整条轨道。本章从积分方程与逐次逼近建立局部存在唯一性，解释局部 Lipschitz 条件的作用，并以 Gronwall 估计连接唯一性、连续依赖与全局控制。最大解和延拓准则区分有限时间爆破与碰到定义域边界；线性系统部分建立矩阵指数、基本解矩阵与变常数公式；流及变分方程说明轨道怎样随初值和参数变化。最后两节选读分别用 Euler 折线证明连续场的 Peano 存在定理，并用 Lyapunov 函数研究平衡点的稳定性。
\end{chapterabstract}

\begin{learninggoals}
\begin{itemize}
\item\label{b2:item:ode-goal-local}能说明非自治方程对时间连续、对状态局部一致 Lipschitz 的具体含义，构造共同短时间的解，并用 Gronwall 估计证明唯一性。
\item\label{b2:item:ode-goal-continuation}能拼接最大解，使用紧集延拓、连续依赖与标量比较，区别局部存在、正向存在和全时间存在。
\item\label{b2:item:ode-goal-linear}能构造连续系数系统的基本解矩阵，计算简单矩阵指数，用变常数公式分析外力，并辨认非交换系数造成的障碍。
\item\label{b2:item:ode-goal-flow}能证明流的复合关系和光滑依赖，写出初值与参数的变分方程，区分自治流与非自治的两参数演化。
\item\label{b2:item:ode-goal-peano}选读后能从 Euler 折线的紧性得到 Peano 存在定理，并解释连续性为何不能保证唯一性。
\item\label{b2:item:ode-goal-stability}选读后能准确区分稳定、渐近稳定与指数稳定，构造简单 Lyapunov 函数，并核验其正定性与沿轨道导数的符号。
\end{itemize}
\end{learninggoals}

速度场给出每一点的运动方向，却不直接给出从某点出发的完整轨道。把微分方程写成积分方程以后，反复代入可以构造轨道；要让轨道决定单值的运动，还须证明唯一性。即使这两步在短时间内成立，轨道仍可能在有限时间内逃向无穷。若要用轨道作局部坐标，还需研究它对初始位置的导数。

\ProseParagraphBreak

\secref{b2:sec:ode-local-existence}从积分方程入手，先建立 Gronwall 估计，再以逐次逼近构造一般非自治局部解；自治 \(C^1\) 情形是后文特征线法最常调用的版本。\secref{b2:sec:ode-continuation}解释为什么短时间结论能够沿轨道继续，以及延拓失败会发生什么。

\ProseParagraphBreak

\secref{b2:sec:ode-linear-systems}既提供显式计算的模型，也为\secref{b2:sec:ode-smooth-flow}的变分方程准备线性工具。矩阵指数由幂级数构造，使标量指数解能够推广到相互耦合的线性系统。初次阅读时应把“解关于时间可微”与“解关于初值可微”分开理解：前者是满足方程的要求，后者需要另做余项估计。

\ProseParagraphBreak

选读的\secref{b2:sec:ode-peano}撤去 Lipschitz 条件，展示“以紧性换取存在”的另一条路线；\secref{b2:sec:ode-lyapunov}把目光从有限时间的构造移到长时间行为。章末前六题覆盖前四节，后两题分别连接 Peano 选解与渐近稳定的模型。

\ProseParagraphBreak

John K. Hunter 与 Bruno Nachtergaele 的《Applied Analysis》\cite[Section 3.5]{Hunter2001Applied}用积分迭代构造局部解，并说明局部存在与全局存在的区别。Gerald Teschl 的《Ordinary Differential Equations and Dynamical Systems》\cite[Section 2.7; Sections 6.5--6.6]{Teschl2012ODE}分别讨论 Euler 近似与 Peano 存在，以及平衡点和 Lyapunov 方法；阅读时可对照紧性证明与唯一性证明各自使用了哪些条件。
% 实读：Hunter--Nachtergaele 作者稿 Section 3.5，印刷73--76页。
% 实读：Teschl 经AMS许可的作者稿 Section 2.7, pp.54--56;
% Sections 6.5--6.6, pp.198, 200--202；引文著录2012正式出版版。

% !TeX root = ../../Book2.tex
\section{初值问题与局部存在唯一}
\label{b2:sec:ode-local-existence}

微分方程规定的是瞬时速度，而我们要寻找的是一整条运动轨道。即使每一点的速度已经给定，也还要回答三个不同的问题：经过指定点的轨道是否存在，是否只有一条，以及它至少能走多长时间。本节用积分方程构造短时间解，并把保证存在的速度界与保证唯一的 Lipschitz 界区分开来。

\subsection{积分方程与 Gronwall 估计}
\label{b2:subsec:ode-integral-gronwall}

\begin{definition}\label{b2:def:ode-initial-value}
设 \(d\geq1\)，\(\Omega\subset\mathbb R\times\mathbb R^d\) 为开集，\(f\colon\Omega\to\mathbb R^d\) 连续，\((t_0,a)\in\Omega\)。寻找满足
\[
 X'(t)=f(t,X(t)),\quad X(t_0)=a
\]
的函数，称为一个初值问题。它在含 \(t_0\) 的非退化区间 \(I\) 上的解，是一个 \(C^1\) 映射 \(X\colon I\to\mathbb R^d\)，要求 \((t,X(t))\in\Omega\)，并逐点满足上述方程；若 \(I\) 含有端点，在端点采用单侧导数。若 \(f(t,x)=b(x)\) 不显含 \(t\)，则方程称为\term[zizhi-fangcheng]{自治方程}[autonomous equation]；其余情形称为\term[feizizhi-fangcheng]{非自治方程}[nonautonomous equation]。
\end{definition}

这里的一个未知函数可以同时有 \(d\) 个分量。二阶标量方程也属于这种框架：令 \(x_1=u\)、\(x_2=u'\)，则 \(u''=F(t,u,u')\) 化为 \(x_1'=x_2\)、\(x_2'=F(t,x_1,x_2)\)。初始位置与初始速度一起组成状态 \((x_1,x_2)\)，不能只给其中一个。

\ProseParagraphBreak
微积分基本定理把初值问题等价地写成
\[
 X(t)=a+\int_{t_0}^t f(s,X(s))\,\mathrm ds.
\]
从微分方程到积分方程只需积分；反过来，只要 \(X\) 连续且满足这个等式，连续的被积函数就使右端属于 \(C^1\)，从而恢复微分方程。这一等价性很有用：构造解时只需在连续函数空间中取极限，无须事先证明近似解的导数收敛。

\ProseParagraphBreak
如果两条曲线代入同一个积分方程，曲线的差会再次出现在积分号内。\term[gronwall-budengshi]{Gronwall 不等式}[Gronwall inequality]把这种含有未知量自身的积分估计变成显式上界。

\begin{lemma}[Gronwall]\label{b2:lem:ode-gronwall}
设 \(h>0\)、\(A,L\geq0\)，连续非负函数 \(q\) 在 \([0,h]\) 上满足
\[
 q(t)\leq A+L\int_0^t q(s)\,\mathrm ds.
\]
则 \(q(t)\leq A e^{Lt}\) 对每个 \(t\in[0,h]\) 成立，包括 \(A=0\) 的情形。
\end{lemma}
\begin{proof}
令 \(Q(t)=A+L\int_0^t q(s)\,\mathrm ds\)。由 \(q\leq Q\) 得 \(Q'\leq LQ\)，从而 \((e^{-Lt}Q)'\leq0\)。因此 \(q(t)\leq Q(t)\leq A e^{Lt}\)。当 \(A=0\) 时，同一微分不等式给出 \(q=0\)，正好对应零初始误差。
\end{proof}

初始误差为零时，结论给出 \(q=0\)，这将证明唯一性；初始误差很小时，指数因子刻画误差在有限时间内可能放大的程度。后一种用途还需要容许方程本身有误差。

\begin{corollary}[带误差项的 Gronwall 估计]\label{b2:cor:ode-gronwall-forcing}
设 \(\alpha<\beta\)，\(A\geq0\)，\(q,k,r\colon[\alpha,\beta]\to[0,\infty)\) 连续，并且
\[
 q(t)\leq A+\int_\alpha^t \bigl(k(s)q(s)+r(s)\bigr)\,\mathrm ds.
\]
则对 \(\alpha\leq t\leq\beta\)，
\[
 q(t)\leq A\exp\!\left(\int_\alpha^t k(s)\,\mathrm ds\right)
 +\int_\alpha^t \exp\!\left(\int_s^t k(\tau)\,\mathrm d\tau\right)r(s)\,\mathrm ds.
\]
特别地，若 \(k\leq L\)、\(r\leq\eta\)，则 \(q(t)\leq\bigl(A+\eta(t-\alpha)\bigr)e^{L(t-\alpha)}\)。
\end{corollary}
\begin{proof}
把假设右端记为 \(Q(t)\)，则 \(q\leq Q\)、\(Q(\alpha)=A\)，且 \(Q'\leq kQ+r\)。乘以积分因子 \(\exp(-\int_\alpha^t k)\) 后积分，得到所述不等式。最后一条估计把两个指数因子都用 \(e^{L(t-\alpha)}\) 控制即可。
\end{proof}

这些估计同样能向初始时刻的左侧使用，但要先改用正向变量 \(s=t_0-t\)。直接把上限小于下限的积分当成非负量，会改变不等号的意义；反向时间的变换避免了这个错误。

\subsection{非自治方程的局部 Lipschitz 条件}
\label{b2:subsec:ode-nonautonomous-lipschitz}

为了比较 \(f(t,x)\) 与 \(f(t,y)\)，需要控制的是状态变量的变化。时间只起参数作用，构造中不要求 \(f\) 对时间可微。不过，控制常数必须在一小段时间内统一，否则无法估计整个积分。

\begin{definition}\label{b2:def:ode-local-lipschitz}
设 \(d\geq1\)，\(\Omega\subset\mathbb R\times\mathbb R^d\) 为开集。称映射 \(f\colon\Omega\to\mathbb R^d\) 对状态变量满足\term[jubu-yizhi-lipschitz]{局部一致 Lipschitz 条件}[locally uniform Lipschitz condition]，是指对每个 \((t_*,x_*)\in\Omega\)，存在包含该点的开柱体 \(J\times B\subseteq\Omega\)，其中 \(J\) 为时间区间、\(B\) 为空间球，以及常数 \(L\geq0\)，使
\[
 |f(t,x)-f(t,y)|\leq L|x-y|
 \quad(t\in J,\ x,y\in B).
\]
“一致”指同一个 \(L\) 适用于这个 \(J\) 中的全部时间。
\end{definition}

若 \(f\) 连续，且对 \(x\) 的偏导数组成的矩阵 \(D_xf\) 存在并关于 \((t,x)\) 连续，就满足这个条件：在一个闭柱体上取 \(\|D_xf\|\) 的上界，再沿连接 \(x,y\) 的线段积分。这里要把状态邻域取成球，保证线段留在求导区域内。反过来，局部一致 Lipschitz 不要求可微，例如实标量函数 \(f(t,x)=|x|+\sin t\) 在 \(x=0\) 不可微，却对 \(x\) 满足常数为 \(1\) 的 Lipschitz 估计。

\ProseParagraphBreak
下面证明\term[picard-lindelof]{Picard--Lindelöf 定理}[Picard--Lindelöf theorem]。定理特意同时允许初始时刻和初始位置在一个小邻域内变化；这种共同存在时间将在拼接解和研究连续依赖时使用。

\begin{theorem}[Picard--Lindelöf]\label{b2:thm:ode-nonautonomous-picard}
设 \(d\geq1\)，\(\Omega\subset\mathbb R\times\mathbb R^d\) 开，\(f\colon\Omega\to\mathbb R^d\) 连续且对状态变量满足局部一致 Lipschitz 条件。给定 \((t_0,a_0)\in\Omega\)，可取 \(\delta,R>0\)，使闭柱体
\[
 Q=[t_0-\delta,t_0+\delta]\times\overline B_R(a_0)
 \subseteq\Omega
\]
上有 \(|f|\leq M\)，并且对同一时刻的任意 \(x,y\in\overline B_R(a_0)\) 有 \(|f(t,x)-f(t,y)|\leq L|x-y|\)。取 \(h>0\) 满足
\[
 h\leq\delta/2,\quad Mh\leq R/2,\quad Lh<1.
\]
则每个满足 \(|\tau-t_0|\leq\delta/2\)、\(|a-a_0|<R/2\) 的初值 \(X(\tau)=a\)，都有一个定义于 \([\tau-h,\tau+h]\) 的解，且解取值于 \(\overline B_R(a_0)\)。任何两个同初值的解在其共同定义区间上相等。
\end{theorem}
\begin{proof}
由局部一致 Lipschitz 条件取一个开柱体，再缩小为包含 \((t_0,a_0)\) 的闭柱体 \(Q\)。连续性和紧性给出有限的速度界 \(M\)。当 \(M\) 或 \(L\) 为零时，相应的时间限制自动成立，所以总能选择正的 \(h\)。

固定 \(\tau,a\)，令 \(E_a\) 为所有满足
\[
 Y\in C([\tau-h,\tau+h];\mathbb R^d),
 \quad \sup_{|t-\tau|\leq h}|Y(t)-a|\leq R/2
\]
的曲线所成的集合。它在一致范数下完备：一致 Cauchy 列逐点有极限，且一致收敛；极限连续，并仍满足闭球约束。对 \(Y\in E_a\) 定义
\[
 (TY)(t)=a+\int_\tau^t f(s,Y(s))\,\mathrm ds.
\]
所有积分都在 \(Q\) 内进行。由 \(|TY(t)-a|\leq Mh\leq R/2\) 知 \(T(E_a)\subseteq E_a\)，而 Lipschitz 条件给出
\[
 \|TY-TZ\|_\infty\leq Lh\,\|Y-Z\|_\infty.
\]
从常曲线 \(Y_0=a\) 开始，依次令 \(Y_{n+1}=TY_n\)。记 \(\theta=Lh<1\)，则
\[
 \|Y_{n+1}-Y_n\|_\infty
 \leq\theta^n\|Y_1-Y_0\|_\infty.
\]
把相邻差求和，得到对 \(m>n\) 的一致估计
\[
 \|Y_m-Y_n\|_\infty
 \leq\frac{\theta^n}{1-\theta}\|Y_1-Y_0\|_\infty.
\]
因此 \(Y_n\) 在 \(E_a\) 中收敛到某个 \(X\)。算子 \(T\) 的 Lipschitz 连续性使 \(TX=X\)，积分方程随即给出 \(X\in C^1\) 及所需初值问题。若 \(X,Z\) 都是不动点，则 \(\|X-Z\|_\infty\leq\theta\|X-Z\|_\infty\)，故它们相等。

还须说明，这种唯一性不局限于所选的曲线集合。任意两条同初值解，由连续性先在一个足够短的区间上留在同一 Lipschitz 柱体内。积分方程作差并使用 \(A=0\) 的 Gronwall 不等式，得它们在该短区间上相等，负时间经反向变量处理。在共同定义区间中，两解相等的时刻集合是非空闭集；在每个相等时刻重新使用刚才的局部论证，又知该集合是相对开集。区间的连通性使它充满整个共同区间。
\end{proof}

证明中两种时间限制承担不同任务：\(Mh\leq R/2\) 防止曲线跑出已知方程的区域，\(Lh<1\) 使一次积分缩小两条曲线的距离。逐次逼近的第 \(n+1\) 条曲线使用第 \(n\) 条曲线上的速度，它不是把未知解代入右端后的循环定义。

\ProseParagraphBreak
可以直接跟算 \(X'=X\)、\(X(0)=a\) 的迭代：从 \(Y_0=a\) 得到 \(Y_1=a(1+t)\)、\(Y_2=a(1+t+t^2/2)\)，一般有 \(Y_n=a\sum_{j=0}^n t^j/j!\)。极限是熟悉的 \(ae^t\)。一般方程未必有闭式表达式，但一致收敛和积分方程仍给出确切的解；存在性并不以能写出初等函数公式为前提。

\subsection{自治 \texorpdfstring{\(C^1\)}{C¹} 方程与共同存在时间}
\label{b2:subsec:ode-autonomous-picard}

后文的特征线和流通常由一个 \(C^1\) 向量场产生。把一般定理专门写成这一形式，可以直接看见邻近初始点共用哪些常数。

\begin{theorem}[Picard--Lindelöf，自治情形]\label{b2:thm:07-local-ode}
设 \(d\geq1\)，\(U\subset\mathbb R^d\) 为开集，\(b\in C^1(U;\mathbb R^d)\)。对每个 \(a_0\in U\)，存在 \(r,h>0\)，使每个 \(a\in B_r(a_0)\) 都有唯一的局部解
\[
 X'(t)=b(X(t)),\quad X(0)=a,\quad |t|\leq h.
\]
各解取值于同一个紧球 \(\overline B_{2r}(a_0)\subseteq U\)。若 \(L\) 是该球上 \(\|Db\|\) 的上界，则对任意 \(a,\widetilde a\in B_r(a_0)\) 及 \(|t|\leq h\)，
\begin{equation}\label{b2:eq:07-flow-lipschitz}
 |X(t,a)-X(t,\widetilde a)|
 \leq e^{L|t|}|a-\widetilde a|.
\end{equation}
局部唯一性指任意两条满足同一初值的解在共同定义区间上相等。
\end{theorem}
\begin{proof}
取闭球 \(\overline B_{2r}(a_0)\subseteq U\)，在其上用 \(M,L\) 分别控制 \(|b|\)、\(\|Db\|\)。沿球内线段积分，得 \(|b(x)-b(y)|\leq L|x-y|\)。对 \(f(t,x)=b(x)\) 使用\cref{b2:thm:ode-nonautonomous-picard}，取 \(Mh\leq r\)、\(Lh<1\)，便得到全部初始点共用的时间 \([-h,h]\)。

两个解的积分方程作差，在正时间得到
\[
 |X(t,a)-X(t,\widetilde a)|
 \leq |a-\widetilde a|
 +L\int_0^t|X(s,a)-X(s,\widetilde a)|\,\mathrm ds.
\]
\cref{b2:lem:ode-gronwall}给出所需估计。负时间的曲线 \(s\mapsto X(-s,a)\) 满足速度场为 \(-b\) 的方程，其 Lipschitz 常数仍为 \(L\)，故得到带 \(|t|\) 的统一表达式。共同区间上的唯一性已包含在一般定理中。
\end{proof}

这里的 \(C^1\) 条件是一个容易核验的充分条件，其作用是提供局部 Lipschitz 常数；一般定理并不要求状态变量处处可微。也不能把连续性误当作唯一性的充分条件。实方程 \(X'=2\sqrt{|X|}\)、\(X(0)=0\) 的零解与下列每条曲线都满足同一初值问题：对任意 \(c\geq0\)，令
\[
 X_c(t)=
 \begin{cases}
 0,&t\leq c,\\
 (t-c)^2,&t\geq c.
 \end{cases}
\]
两段在 \(t=c\) 的导数均为零，因此它确实是 \(C^1\) 解。轨道可以在原点等待任意时间，再开始运动；\(2\sqrt{|x|}\) 在原点附近的差商无界，恰好不满足局部 Lipschitz 条件。仅有连续性时如何保留存在性，将在\secref{b2:sec:ode-peano}中另行证明。

% !TeX root = ../../Book2.tex
\section{最大解、延拓与连续依赖}
\label{b2:sec:ode-continuation}

短时间构造不规定解的最终时间范围。唯一性使不同局部构造能够相互接合；增长估计则用于判断这种接合能否经过任意有限时间。本节先在自治 \(C^1\) 条件下整理这两个步骤。

\subsection{最大解区间}
\label{b2:subsec:ode-maximal-interval}

设 \(U\subset\mathbb R^d\) 开，\(b\in C^1(U;\mathbb R^d)\)、\(a\in U\)。考虑所有在含 \(0\) 的开区间上满足 \(X'=b(X)\)、\(X(0)=a\) 的解。它们在重叠区间上由\cref{b2:thm:07-local-ode}的唯一性相等，因此可以拼成一条解。所有这些区间的并仍为含 \(0\) 的开区间，记为 \(I_a\)。每个时刻附近的拼接函数都等于某条局部解，所以它属于 \(C^1\) 并满足方程。这就得到经过 \(a\) 的最大解；任何同初值解都是它的限制。

\subsection{共同短区间上的初始点依赖}
\label{b2:subsec:ode-initial-dependence}

在\cref{b2:thm:07-local-ode}选定的共同紧球和共同时间 \([-h,h]\) 上，式\eqref{b2:eq:07-flow-lipschitz}给出
\[
 \sup_{|t|\leq h}|X(t,a)-X(t,\widetilde a)|
 \leq e^{Lh}\cdot|a-\widetilde a|.
\]
因此解在一致范数下连续地依赖初始点，且此处具有 Lipschitz 控制。这个共同时间来自局部球上的速度界和导数界，不能直接当作所有初始点的共同全局时间。

\subsection{线性增长与全时间存在}
\label{b2:subsec:ode-linear-growth}

若速度至多随位置线性增长，Gronwall 估计可以把任意有限时间内的轨道限制在紧球中。

\begin{proposition}\label{b2:prop:ode-linear-growth}
设 \(d\) 为正整数，\(b\in C^1(\mathbb R^d;\mathbb R^d)\)，且存在 \(C\geq0\) 使对所有 \(x\in\mathbb R^d\) 都有
\[
 |b(x)|\leq C(1+|x|).
\]
则对每个 \(a\in\mathbb R^d\)，初值问题 \(X'=b(X)\)、\(X(0)=a\) 的最大解在整个实轴上存在，并满足
\[
 1+|X(t,a)|\leq(1+|a|)\cdot e^{C|t|}.
\]
\end{proposition}
\begin{proof}
对正时间，积分方程给出
\[
 1+|X(t,a)|\leq1+|a|+C\int_0^t(1+|X(s,a)|)\,\mathrm ds.
\]
在最大解区间内的任意紧时间区间上应用\cref{b2:lem:ode-gronwall}，得到所述估计；负时间对反向速度 \(-b\) 使用同一论证。若最大解区间有有限右端点 \(T\)，此估计使轨道在 \([0,T)\) 落于某个紧球。该球上速度有界，所以 \(|X(t,a)-X(s,a)|\leq M|t-s|\)，从而 \(t\uparrow T\) 时有极限 \(a_T\)。由 \(a_T\) 出发的局部解可与原解拼接；积分方程保证拼接点导数也满足方程，与最大性矛盾。左端点相同。
\end{proof}

线性增长条件不能无条件省去。对于 \(b(x)=x^2\)，从 \(a>0\) 出发的轨道为 \(a/(1-at)\)，在 \(t=1/a\) 逃向无穷。局部唯一性仍成立，但这条最大解不能延续到所有正时间。另一方面，若只关心一条有限轨道及其附近，紧集上的局部条件已经足够，不必强求整个空间上的有界导数。

\subsection{一般延拓准则}
\label{b2:subsec:ode-general-continuation}

待撰写。拟在一般定义域中区分离开每个紧集、到达定义域边界和状态无界等情形，并补充非自治问题的延拓与连续依赖。

\subsection{标量比较}
\label{b2:subsec:070207}

待撰写。

% !TeX root = ../../Book2.tex
\section{线性系统、矩阵指数与变常数}
\label{b2:sec:ode-linear-systems}

线性系统的时间系数即使只有连续性，也能由积分迭代构造唯一解。这里的估计含有阶乘，不必把时间区间缩短到使积分算子成为压缩；后文 Sturm--Liouville 问题的通量系统正需要这个版本。矩阵指数与一般变常数公式仍留待撰写。

\subsection{连续系数线性系统与积分迭代}
\label{b2:subsec:070205}

\begin{proposition}\label{b2:prop:ode-continuous-linear}
设 \(n\) 为正整数，\(\mathbb K=\mathbb R\) 或 \(\mathbb C\)，\(\alpha<\beta\)，\(A\in C([\alpha,\beta];\mathbb K^{n\times n})\)。对任意 \(t_0\in[\alpha,\beta]\)、\(y_0\in\mathbb K^n\)，问题
\[
 Y'(t)=A(t)Y(t),\quad Y(t_0)=y_0
\]
在整个 \([\alpha,\beta]\) 上有唯一的 \(C^1\) 解；端点处的导数取相应单侧值。若 \(M=\max_{[\alpha,\beta]}\|A(t)\|\)，其中矩阵范数由 Euclid 范数诱导，则
\[
 |Y(t)|\leq e^{M|t-t_0|}\cdot|y_0|.
\]
\end{proposition}
\begin{proof}
令 \(Y_0(t)=y_0\)，递推定义
\[
 Y_{k+1}(t)=y_0+\int_{t_0}^t A(s)Y_k(s)\,\mathrm ds,
 \quad k\geq0.
\]
第一次差满足 \(|Y_1(t)-Y_0(t)|\leq M|t-t_0|\cdot|y_0|\)。反复积分并归纳得到
\[
 |Y_{k+1}(t)-Y_k(t)|
 \leq\frac{M^{k+1}|t-t_0|^{k+1}}{(k+1)!}\cdot|y_0|.
\]
当 \(t<t_0\) 时，把积分方向反转后使用同一估计。由于 \(|t-t_0|\leq\beta-\alpha\)，右端的数项级数收敛，故 \(Y_k\) 一致收敛到连续函数 \(Y\)。系数有界，允许在积分方程中取一致极限，得到
\[
 Y(t)=y_0+\int_{t_0}^t A(s)Y(s)\,\mathrm ds.
\]
被积函数连续，所以 \(Y\in C^1\) 并满足初值问题。把逐次差的估计求和，同时得到 \(|Y(t)|\leq e^{M|t-t_0|}\cdot|y_0|\)。

若 \(Y,\widetilde Y\) 是同初值的两解，令 \(Z=Y-\widetilde Y\)、\(D=\max_{[\alpha,\beta]}|Z|\)。由 \(Z(t)=\int_{t_0}^t A(s)Z(s)\,\mathrm ds\) 反复代入，对每个正整数 \(m\) 都有
\[
 |Z(t)|\leq\frac{M^m|t-t_0|^m}{m!}\cdot D.
\]
令 \(m\to\infty\) 得 \(Z(t)=0\)，故解唯一。
\end{proof}

矩阵初值问题可逐列应用本命题。若指定某时刻的初始矩阵，每一列唯一确定，因而整个矩阵解也唯一确定。证明只使用系数连续与紧区间上的有界性，没有对 \(A\) 求导。

\subsection{矩阵指数与变常数公式}
\label{b2:subsec:070206}

待撰写。拟先处理常系数矩阵指数，再建立基本解矩阵及非齐次系统的变常数公式。前一命题目前只给出连续系数齐次系统的初值理论。

% !TeX root = ../../Book2.tex
\section{光滑依赖与流}
\label{b2:sec:ode-smooth-flow}

唯一性使从不同时间出发的轨道能够相互接合，光滑依赖则允许把轨道用于坐标变换。本节在自治 \(C^1\) 条件下建立这两项性质；高阶光滑依赖尚待撰写。

\subsection{流映射与复合关系}
\label{b2:subsec:070202}

设 \(U\subset\mathbb R^d\) 开，\(b\in C^1(U;\mathbb R^d)\)、\(a\in U\)。按\secref{b2:sec:ode-continuation}的拼接构造，初值问题 \(X'=b(X)\)、\(X(0)=a\) 有最大解区间 \(I_a\)。把从 \(a\) 出发的最大解记为 \(\Phi_t(a)\)，称 \(a\mapsto\Phi_t(a)\) 为\term[liu-yingshe]{流映射}[flow map]。记号 \(\Phi_t(a)\) 只在 \(t\in I_a\) 时使用；给映射命名并不意味着它对所有时间存在。
\symidx{\(\Phi_t\)：自治速度场的流映射}

\begin{proposition}\label{b2:prop:07-flow-composition}
设 \(d\) 为正整数，\(U\subset\mathbb R^d\) 为开集，\(b\in C^1(U;\mathbb R^d)\)，\(\Phi_t(a)\) 是 \(X'=b(X)\)、\(X(0)=a\) 的最大局部流。凡各项有定义，流满足
\[
 \Phi_0(a)=a,\quad
 \Phi_{t+s}(a)=\Phi_t\bigl(\Phi_s(a)\bigr),\quad
 \Phi_{-t}\bigl(\Phi_t(a)\bigr)=a.
\]
固定 \(t\) 时，\(\Phi_t\) 在其定义域上单射，其逆为相应的 \(\Phi_{-t}\)。
\end{proposition}
\begin{proof}
曲线 \(r\mapsto\Phi_{s+r}(a)\) 与 \(r\mapsto\Phi_r(\Phi_s(a))\) 在 \(r=0\) 具有相同初值，且都满足自治方程。唯一性给出复合关系。令 \(r=-s\) 得逆关系，进而得到单射性。
\end{proof}

\subsection{一阶变分方程}
\label{b2:subsec:ode-first-variation}

\begin{theorem}\label{b2:thm:07-flow-derivative}
设 \(d\) 为正整数，\(U\subset\mathbb R^d\) 为开集，\(b\in C^1(U;\mathbb R^d)\)，\(\Phi_t(a)\) 是 \(X'=b(X)\)、\(X(0)=a\) 的最大局部流。流 \((t,a)\mapsto\Phi_t(a)\) 在其开定义域上属于 \(C^1\)。对初始点的导数 \(A(t,a)=D_a\Phi_t(a)\) 满足
\begin{equation}\label{b2:eq:07-variational-ode}
 A'(t,a)=Db\bigl(\Phi_t(a)\bigr)A(t,a),\quad A(0,a)=I.
\end{equation}
因此固定时间的流是其定义域到像集的 \(C^1\) 微分同胚。
\end{theorem}
\begin{proof}
先在\cref{b2:thm:07-local-ode}构造的紧球和短时间内工作。式\eqref{b2:eq:07-flow-lipschitz}与 \(|\partial_t\Phi_t|\leq M\) 保证流联合连续。矩阵积分方程
\[
 A(t)=I+\int_0^t Db\bigl(\Phi_s(a)\bigr)A(s)\,\mathrm ds
\]
的系数连续，逐列应用\cref{b2:prop:ode-continuous-linear}即可构造解；其范数的积分不等式与\cref{b2:lem:ode-gronwall}给出 \(\|A(t)\|\leq e^{L|t|}\)。现取小向量 \(v\)，令 \(\delta_v(t)=\Phi_t(a+v)-\Phi_t(a)\)。中值积分给出
\[
 b\bigl(\Phi_t(a+v)\bigr)-b\bigl(\Phi_t(a)\bigr)
 =Db\bigl(\Phi_t(a)\bigr)\delta_v(t)+r_v(t).
\]
导数 \(Db\) 在紧球上一致连续。记其连续模为 \(\omega(\eta)\)，即距离至多 \(\eta\) 的两点上导数差的上确界；则 \(\omega(\eta)\to0\)。由 \(|\delta_v(t)|\leq e^{Lh}\cdot|v|\)，中值积分的余项满足
\[
 \sup_{|t|\leq h}|r_v(t)|
 \leq e^{Lh}\cdot|v|\cdot\omega(e^{Lh}\cdot|v|).
\]
令 \(e_v=\delta_v-Av\)，相减后的积分方程给出
\[
 |e_v(t)|\leq L\int_0^{|t|}|e_v(\pm s)|\,\mathrm ds
              +h\sup_{|s|\leq h}|r_v(s)|,
\]
其中正负号与 \(t\) 的符号一致。Gronwall 估计因此给出
\[
 \sup_{|t|\leq h}|e_v(t)|
 \leq h e^{2Lh}\cdot|v|\cdot\omega(e^{Lh}\cdot|v|).
\]
除以 \(|v|\) 并令 \(v\to0\)，得到
\[
 \sup_{|t|\leq h}
 \frac{|\Phi_t(a+v)-\Phi_t(a)-A(t,a)v|}{|v|}\longrightarrow0.
\]
这证明的是完整 Fréchet 微分，而非只沿坐标求导。若 \(a\) 改为邻近点，矩阵方程的系数一致收敛；作差和 Gronwall 估计又给出 \(A\) 对 \((t,a)\) 连续。时间导数 \(\partial_t\Phi_t(a)=b(\Phi_t(a))\) 也连续，因此流为 \(C^1\)。

对一段已存在的轨道，其图像在任意紧时间区间内为紧集。用上述局部构造覆盖它，并沿轨道取有限分段，连续依赖使邻近初始点的解也能逐段通过这些邻域。故流的定义域为开集，局部可微结论可以沿轨道继续。最后，\cref{b2:prop:07-flow-composition}给出双方均 \(C^1\) 的逆映射，所以导数矩阵可逆，像集也为开集。
\end{proof}

流的空间可微性比轨道本身可微更强。前者控制相邻粒子之间的距离、角度和体积，后者只说明单个粒子的速度存在。后面特征线法中 Jacobi 行列式的演化式依赖的是\eqref{b2:eq:07-variational-ode}，不能只从每条轨道 \(C^1\) 便直接写下。

\subsection{高阶光滑依赖}
\label{b2:subsec:ode-higher-dependence}

待撰写。本节目前只证明自治 \(C^1\) 速度场的一阶光滑依赖。

% !TeX root = ../../Book2.tex
\OptionalSection{连续场与 Peano 存在}{b2:sec:ode-peano}

局部 Lipschitz 条件在前面的证明中承担了两项工作：它使积分迭代能够收敛，也使两个具有相同初值的解不能分开。第二项工作不能只靠连续性完成；第一项却有另一条道路。我们先作出一列容易计算的近似轨道，再从中取出一致收敛的子列，最后证明极限满足积分方程。这里产生收敛的原因是有限维紧性。

\subsection{冻结速度与共同存在区间}
\label{b2:subsec:ode-euler-polygons}

设 \(\Omega\subset\mathbb R\times\mathbb R^d\) 开，\(f:\Omega\to\mathbb R^d\) 连续，\((t_0,a)\in\Omega\)。若真解已经存在，则在很短的时间内有
\[
 x(t_0+\Delta)=a+\Delta f(t_0,a)+o(\Delta).
\]
忽略余项，就得到一步近似：先按当前时刻和当前位置的速度走一小段，再在新位置重新计算速度。取 \(T,r>0\)，使紧柱体
\[
 K=[t_0-T,t_0+T]\times\overline{B_r(a)}\subset\Omega,
 \quad M=\max_K|f|.
\]
先设 \(M>0\)，选 \(0<h\leq\min\{T,r/M\}\)。把 \([t_0,t_0+h]\) 等分成 \(N\) 段，记 \(\Delta_N=h/N\)、\(t_j=t_0+j\Delta_N\)。令 \(x_N(t_0)=a\)，逐段规定
\begin{equation}\label{b2:eq:ode-euler-polygon}
 x_N(t)=x_N(t_j)+(t-t_j)f(t_j,x_N(t_j)),
 \quad t_j\leq t\leq t_{j+1}.
\end{equation}
所得连续分段仿射曲线称为\term[euler-zhexian]{Euler 折线}[Euler polygonal approximation]；节点上的递推就是\term[euler-fangfa]{Euler 方法}[Euler method]。节点处不要求左右导数相等，因此折线本身通常还不是微分方程的解。

\ProseParagraphBreak

对于 \(x'=x\)、\(x(0)=1\)，取步长 \(\Delta=0.3\)，依次得到节点值 \(1,1.3,1.69,2.197\)。\Cref{b2:fig:ode-euler-polygon}画出这三步：每一段采用该段左端点处的速度。图中的折线低于真解，是这个具体算例的性质；一般证明只使用共同有界性和余差趋零。

\begin{figure}[htbp]
\centering
\begin{tikzpicture}[x=7.2cm,y=1.05cm,font=\small]
 \draw[->] (-0.03,0)--(1.03,0) node[right] {\(t\)};
 \draw[->] (0,-0.1)--(0,2.75) node[above] {\(x\)};
 \draw[AnalysisBlue,thick,domain=0:0.94,samples=60]
   plot (\x,{exp(\x)});
 \node[AnalysisBlue,right] at (0.94,2.57) {\(e^t\)};
 \draw[thick] (0,1)--(0.3,1.3)--(0.6,1.69)--(0.9,2.197);
 \node[right] at (0.9,2.12) {Euler 折线};
 \foreach \t/\z in {0/1,0.3/1.3,0.6/1.69,0.9/2.197} {
   \fill (\t,\z) circle (1.5pt);
 }
 \foreach \t in {0.3,0.6,0.9} {
   \draw (\t,0.04)--(\t,-0.04) node[below] {\(\t\)};
 }
 \node[below left] at (0,0) {\(0\)};
 \node[left] at (0,1) {\(1\)};
 \draw[densely dashed,gray] (0.3,0)--(0.3,1.3);
 \draw[densely dashed,gray] (0.6,0)--(0.6,1.69);
 \draw[densely dashed,gray] (0.9,0)--(0.9,2.197);
\end{tikzpicture}
\caption{\(x'=x\) 的 Euler 折线：在每个网格段冻结左端点速度}
\label{b2:fig:ode-euler-polygon}
\end{figure}

先验证每一段上的求值位置始终留在速度场受控的区域内。第一段的起点在 \(K\) 内，故斜率的模不超过 \(M\)。若已经构造到 \(t_j\)，则下一段每个点满足
\[
 |x_N(t)-a|\leq M(t-t_0)\leq Mh\leq r.
\]
归纳可知，每一步求值的位置都在 \(K\) 内，整条折线都可以构造出来。这里的共同时间 \(h\) 与网格数 \(N\) 无关。把任意两个时刻之间的折线段长度相加，还得到
\begin{equation}\label{b2:eq:ode-euler-uniform-bounds}
 |x_N(t)-a|\leq r,
 \quad |x_N(t)-x_N(s)|\leq M|t-s|.
\end{equation}
第一条估计限制轨道的位置，第二条阻止轨道在越来越短的时间内发生固定大小的跳动。两条合起来才足以产生下面的一致收敛子列。

\subsection{从有限维有界性取得一致收敛}
\label{b2:subsec:ode-euler-compactness}

一致有界本身并不保证函数列有一致收敛子列。例如 \(t\mapsto\sin Nt\) 的振荡随 \(N\) 增快。Euler 折线还满足同一个连续性估计，这个额外条件排除了这种现象。下面证明\term[arzela-ascoli-dingli]{Arzel\`a--Ascoli 定理}[Arzel\`a--Ascoli theorem]在有限维值连续函数上的一个版本。

\begin{lemma}[Arzel\`a--Ascoli 的区间版本]\label{b2:lem:ode-arzela-ascoli}
设 \(I=[\alpha,\beta]\subset\mathbb R\) 为非空紧区间，\(d\geq1\)，\(u_n\in C(I;\mathbb R^d)\)。假设存在 \(C<\infty\) 使 \(|u_n(t)|\leq C\) 对所有 \(n,t\) 成立；又假设对每个 \(\varepsilon>0\)，存在与 \(n\) 无关的 \(\delta>0\)，使
\[
 |t-s|<\delta\ \Longrightarrow\ |u_n(t)-u_n(s)|<\varepsilon
 \quad(n\geq1).
\]
则 \((u_n)\) 有在 \(I\) 上一致收敛的子列，其极限连续。
\end{lemma}
\begin{proof}
选 \(I\) 中的可数稠密集 \(D=\{s_1,s_2,\ldots\}\)。在 \(s_1\) 处，序列 \((u_n(s_1))\) 是 \(\mathbb R^d\) 中的有界列，因而有收敛子列。在这个子列中再抽取在 \(s_2\) 处收敛的子列，如此继续。取第 \(j\) 次所选子列的第 \(j\) 项，得到一列 \((v_j)\)：对每个固定的 \(s_k\)，从第 \(k\) 项以后都属于第 \(k\) 次的子列，故 \((v_j(s_k))\) 收敛。

还需要把稠密集上的逐点收敛变成整个区间上的一致收敛。给定 \(\varepsilon>0\)，按共同连续性条件选 \(\delta>0\)，使时间距离小于 \(\delta\) 时函数值之差小于 \(\varepsilon/3\)。利用稠密性与 \(I\) 的紧性，可选有限多个点 \(s_{k_1},\ldots,s_{k_m}\)，使每个 \(t\in I\) 都距其中某点小于 \(\delta\)。这有限多个点上的收敛给出 \(J\)，使 \(j,l\geq J\) 时有
\[
 |v_j(s_{k_i})-v_l(s_{k_i})|<\varepsilon/3
 \quad(1\leq i\leq m).
\]
对于任意 \(t\in I\)，选相应的 \(s_{k_i}\)，得到
\[
 \begin{aligned}
 |v_j(t)-v_l(t)|
 &\leq |v_j(t)-v_j(s_{k_i})|
       +|v_j(s_{k_i})-v_l(s_{k_i})|\\
 &\quad+|v_l(s_{k_i})-v_l(t)|<\varepsilon.
 \end{aligned}
\]
因此 \((v_j)\) 是一致 Cauchy 列。在每个 \(t\) 处取 \(\mathbb R^d\) 中的极限，以上估计说明收敛是一致的。给定一点，把极限函数与某个 \(v_j\) 的两项一致误差、以及 \(v_j\) 在该点的连续性误差分别控制，就证明极限连续。
\end{proof}

引理中的共同连续性条件称为\term[dengdu-lianxu]{等度连续性}[equicontinuity]。证明清楚显示了有限维假设的使用位置：固定时刻的有界向量列必须能取出收敛子列。一般无限维 Banach 空间的有界集未必具有这种性质，不能只把 \(\mathbb R^d\) 换成任意 Banach 空间就照搬本节的存在性证明。

\subsection{Peano 存在定理与非唯一性}
\label{b2:subsec:ode-peano-proof}

现在有了收敛子列，但一致收敛并不自动保持导数。我们不会尝试让折线的导数逐点收敛，而是估计它偏离积分方程的大小。所得的局部存在性结论称为\term[peano-cunzai-dingli]{Peano 存在定理}[Peano existence theorem]。

% 实读核对：Teschl, Ordinary Differential Equations and Dynamical Systems,
% AMS 2012 作者授权副本 Section 2.7, pp.54--56, Theorems 2.18--2.19。
% 正文采用独立展开的有限维紧性和积分余差证明；书目由章首统一引入。
\begin{theorem}[Peano]\label{b2:thm:ode-peano}
设 \(d\geq1\)，\(\Omega\subset\mathbb R\times\mathbb R^d\) 开，\(f:\Omega\to\mathbb R^d\) 连续，\((t_0,a)\in\Omega\)。取 \(T,r>0\)，使
\[
 K=[t_0-T,t_0+T]\times\overline{B_r(a)}\subset\Omega,
 \quad M=\max_K|f|.
\]
若 \(M>0\)，令 \(h=\min\{T,r/M\}\)；若 \(M=0\)，令 \(h=T\)。则初值问题
\[
 x'(t)=f(t,x(t)),\quad x(t_0)=a
\]
至少有一个 \(C^1\) 解定义在 \([t_0-h,t_0+h]\) 上，且 \(|x(t)-a|\leq r\)。区间端点的导数按单侧导数理解。
\end{theorem}
\begin{proof}
当 \(M=0\) 时，\(f\) 在 \(K\) 上恒为零，常值函数 \(x(t)=a\) 就是所需的解。以下设 \(M>0\)。在 \([t_0,t_0+h]\) 上按式\eqref{b2:eq:ode-euler-polygon}构造折线 \(x_N\)。式\eqref{b2:eq:ode-euler-uniform-bounds}给出一致有界性和等度连续性。由\cref{b2:lem:ode-arzela-ascoli}，有子列 \(x_{N_k}\) 一致收敛到连续函数 \(x\)；极限仍满足 \(x(t_0)=a\) 以及 \(|x(t)-a|\leq r\)。

对于每个 \(N\)，令 \(\tau_N(s)\) 为包含 \(s\) 的网格段的左端点；在最终端点处如何赋值不影响积分。逐段积分可得
\[
 x_N(t)=a+\int_{t_0}^t
       f(\tau_N(s),x_N(\tau_N(s)))\,\mathrm ds.
\]
由于 \(f\) 在紧集 \(K\) 上一致连续，可以定义其连续模
\[
 \omega(\rho)=\sup\bigl\{|f(t,y)-f(s,z)|:
 (t,y),(s,z)\in K,\ |t-s|+|y-z|\leq\rho\bigr\},
\]
并有 \(\omega(\rho)\to0\) 当 \(\rho\downarrow0\)。这里的上确界有限，因为 \(f\) 在 \(K\) 上有界。对每个 \(s\)，网格宽度和速度界给出
\[
 |s-\tau_N(s)|\leq\Delta_N,\quad
 |x_N(s)-x_N(\tau_N(s))|\leq M\Delta_N.
\]
因此积分方程的余差满足
\[
 \sup_{t_0\leq t\leq t_0+h}
 \left|x_N(t)-a-\int_{t_0}^t f(s,x_N(s))\,\mathrm ds\right|
 \leq h\omega((1+M)\Delta_N)\longrightarrow0.
\]
又因 \(x_{N_k}\to x\) 一致，\(f\) 的同一个一致连续性保证
\[
 \sup_{s\in[t_0,t_0+h]}|f(s,x_{N_k}(s))-f(s,x(s))|
 \leq\omega(\|x_{N_k}-x\|_\infty)\longrightarrow0.
\]
后一误差的积分至多再乘区间长度 \(h\)。在积分方程余差中取极限，得到
\[
 x(t)=a+\int_{t_0}^t f(s,x(s))\,\mathrm ds.
\]
被积函数连续，故微积分基本定理给出 \(x\in C^1\) 以及所需微分方程。

为构造负向解，对 \(0\leq s\leq h\) 考虑
\[
 y'(s)=-f(t_0-s,y(s)),\quad y(0)=a.
\]
它具有相同的紧柱体速度界，故上面的构造适用。令 \(x(t)=y(t_0-t)\) 便得到 \([t_0-h,t_0]\) 上的解。两侧在 \(t_0\) 取值均为 \(a\)，左右导数均为 \(f(t_0,a)\)，所以可以拼成整个闭区间上的 \(C^1\) 解。
\end{proof}

Peano 存在定理所保证的是至少一条轨道。它没有给出压缩映射中的几何误差估计，也没有说明不同抽取方式会得到同一极限。若另外已有局部 Lipschitz 条件，唯一性便会使每个收敛子列的极限相同；若整个折线列还不一致收敛，就能选出始终与唯一解相距至少某个正数的子列，再由紧性从中抽出收敛子列，产生矛盾。因此在有唯一性的情形，紧性论证还可以升级为整个折线列的收敛。

\ProseParagraphBreak

回看\secref{b2:sec:ode-local-existence}已经验证的初值问题 \(x'=2\sqrt{|x|}\)、\(x(0)=0\)。除了始终停在零点的解，还有等待到任意 \(c\geq0\) 再沿 \((t-c)^2\) 运动的解。Peano 定理适用于这个连续场，保证存在，却没有提供决定何时离开的额外条件。

\ProseParagraphBreak

唯一性的缺失也妨碍把这些轨道自然组织成前一节的可微流。近似算法或初值扰动可能选出不同轨道，\cref{b2:ex:ode-peano-selection}给出一个具体计算。后续特征线法所需的可逆流仍应使用前四节的存在唯一与光滑依赖结果，不能仅以连续场存在性替换。\Cref{b2:tab:ode-regularity}集中比较本章已经证明的三种局部结论；任何一行都不能单凭所列正则性推出全时间存在。

\begin{table}[htbp]
\centering
\small
\caption[速度场条件与局部解的性质]{速度场条件与局部解的性质。状态空间有限维，方程定义于开集；Lipschitz 条件关于时间局部一致，光滑性指全部变量。}
\label{b2:tab:ode-regularity}
\begin{tabularx}{\linewidth}{@{}>{\raggedright\arraybackslash}p{2.8cm}>{\raggedright\arraybackslash}X>{\raggedright\arraybackslash}X@{}}
\toprule
速度场条件 & 得到的局部性质 & 仍需另外核验的性质 \\
\midrule
仅连续 & Peano 保证至少一个 \(C^1\) 解。 & 唯一性、单值的初值依赖。 \\
连续且对状态局部 Lipschitz & 解唯一，有限轨道段上连续依赖初值。 & 对初值的可微性，见\cref{b2:ex:ode-lipschitz-nonsmooth-flow}。 \\
属于 \(C^k\)，\(k\geq1\) & 解的演化关于时间、初值及光滑参数属于 \(C^k\)。 & 最大存在区间及长期稳定性。 \\
\bottomrule
\end{tabularx}
\end{table}

% !TeX root = ../../Book2.tex
\OptionalSection{平衡点与 Lyapunov 稳定性}{b2:sec:ode-lyapunov}

待撰写。拟引入平衡点、稳定性与渐近稳定性，并用 Lyapunov 函数给出基本判据。本节范围限于这些基本概念与方法。


\begin{chaptersummary}
局部存在唯一性把一个微分关系转化为初始数据决定的轨道。积分迭代中的共同时间保证所有候选曲线留在同一个受控区域，Lipschitz 常数则使误差收缩。Gronwall 估计在构造完成后继续发挥作用：它控制相邻初值和相近方程的解，使短时间结论能够覆盖已经存在的有限轨道段。

\ProseParagraphBreak

最大解只是不能再延长的解，并不自动在整个实轴存在。有限端点意味着时空轨道最终离开方程定义域中的每个紧集；在整个状态空间上，这才具体化为状态的无界增长。线性增长、比较或其他先验界能够排除相应的逃逸，使用时要核对所控制的是正向还是双向时间。

\ProseParagraphBreak

连续系数线性系统的积分迭代具有阶乘界，因此不必在每一小段重新压缩时间。基本解矩阵把任意初值的解组织为一个线性映射，变常数公式再累计各时刻外力的影响。常系数时矩阵指数给出这一演化；变系数矩阵不交换时，指数积分公式一般失效。

\ProseParagraphBreak

自治流的复合与逆关系来自唯一性，空间可微性来自变分方程及余项估计。把轨道与其导数放进增广系统，可以继续证明高阶光滑性；再把时间和常参数作为状态变量，便得到非自治演化的联合光滑性。每一步都保留了正则性条件，不能从轨道自身可微直接跨到任意阶参数可微。

\ProseParagraphBreak

选读部分展示两种不同的问题。Peano 定理只要求连续场，但以子列紧性构造的解可能不唯一；Lyapunov 方法则在唯一性框架下，用一个沿轨道下降的函数控制全部正向时间。非正导数通常只给稳定，严格负导数给出渐近稳定，而指数速率还需要定量的二次型控制。
\end{chaptersummary}

\BookExercises

% 目标：通过显式解辨认最大区间、前后时间和定义域内的爆破。
% 来源：自拟；logistic模型的分类计算，正文不依赖本题。
\begin{exercise}[同一方程的不同最大区间]\label{b2:ex:ode-logistic}
对每个 \(a\in\mathbb R\)，求初值问题 \(x'=x(1-x)\)、\(x(0)=a\) 的最大解及其最大时间区间。分别讨论 \(a<0\)、\(a=0\)、\(0<a<1\)、\(a=1\)、\(a>1\)。哪些初值的解对全部正向时间存在？它们是否都在全部反向时间存在？
\end{exercise}
\begin{solution}
\(a=0\) 与 \(a=1\) 分别给出恒等于 \(0\) 与 \(1\) 的全时间解。对其余初值，在不碰到这两个平衡值的区间上分离变量，得到
\[
 \log\left|\frac{x}{1-x}\right|=t+\text{常数},\quad
 x(t)=\frac{ae^t}{1-a+ae^t}.
\]
也可以直接代入验证右式，随后用局部唯一性确定它就是初值解。因此不必先假定分离变量过程已覆盖全部轨道；最大区间由分母的零点决定。

若 \(0<a<1\)，分母对所有实数 \(t\) 都正，故最大区间为 \(\mathbb R\)。若 \(a<0\)，分母唯一零点为 \(T=\log(1-1/a)>0\)，含 \(0\) 的最大区间是 \((-\infty,T)\)，且 \(t\uparrow T\) 时 \(x(t)\to-\infty\)。若 \(a>1\)，相同零点满足 \(T<0\)，最大区间是 \((T,\infty)\)，且 \(t\downarrow T\) 时 \(x(t)\to+\infty\)。无穷极限阻止了穿过该时刻的实值连续延拓。

因此恰好 \(a\geq0\) 的解在全部正向时间存在；其中 \(a>1\) 的解不能在全部反向时间存在。对任意 \(a>0\)，\(t\to+\infty\) 时 \(x(t)\to1\)。多项式速度场在整个实轴光滑，只保证局部存在唯一；它既不能排除其他初值的正向爆破，也不能从正向有界推出双向全局存在。
\end{solution}

% 目标：以比较获得非线性先验界，不把超线性增长误判为必然爆破。
% 来源：自拟；有界外力与耗散项的标量组合。
\begin{exercise}[超线性速度与正向控制]\label{b2:ex:ode-forced-dissipation}
设 \(g\in C(\mathbb R;\mathbb R)\)，\(|g(t)|\leq M\)，其中 \(M\geq0\)。证明 \(x'=-x^3+g(t)\)、\(x(0)=a\in\mathbb R\) 的解在所有 \(t\geq0\) 存在，且
\[
 |x(t)|\leq R\coloneqq\max\{|a|,M^{1/3}\}\quad(t\geq0).
\]
用 \(g=0\)、\(a=1\) 说明这一结论不能直接改成全时间存在。
\end{exercise}
\begin{solution}
速度场连续，状态导数 \(-3x^2\) 在每个紧区间有界，因此局部存在唯一。常函数 \(R\) 满足 \(0\geq-R^3+g(t)\)，是上解；常函数 \(-R\) 满足 \(0\leq R^3+g(t)\)，是下解。标量比较分别用于解与这两个常函数，在解已有的全部正向时间上得到 \(-R\leq x(t)\leq R\)。当 \(R=0\) 时也可直接由唯一性得到零解。

若最大右端点 \(T\) 有限，轨道的时空图像位于紧集 \([0,T]\times[-R,R]\subset\mathbb R^2\)。一般延拓准则排除这一情形，故解存在于所有正向时间。这里三次项指向原点，其绝对值的超线性增长并不引起正向爆破。

当 \(g=0\)、\(a=1\) 时，解为 \(x(t)=(1+2t)^{-1/2}\)，最大区间 \((-1/2,\infty)\)。它在正向时间衰减，却在反向有限时刻爆破。这也说明使用先验界时必须保留时间方向。
\end{solution}

% 目标：从系统变常数计算共振，辨认固定时间连续依赖与长期一致控制。
% 来源：自拟；经典受迫振子模型重新组织。
\begin{exercise}[受迫振子与共振极限]\label{b2:ex:ode-resonance}
设 \(\omega>0\)、\(\nu\geq0\)。把
\[
 y''+\omega^2y=\cos(\nu t),\quad y(0)=y'(0)=0
\]
写成二维一阶系统，并用变常数公式求解。分别给出 \(\nu\ne\omega\) 与 \(\nu=\omega\) 的表达式，说明 \(\nu\to\omega\) 时解在每个紧时间区间一致收敛，却不能由此得到对全部 \(t\geq0\) 通用的有界响应估计。
\end{exercise}
\begin{solution}
取 \(Y=(y,y')^{\mathsf T}\)，则
\[
 Y'=AY+\binom{0}{\cos(\nu t)},\quad
 A=\begin{pmatrix}0&1\\-\omega^2&0\end{pmatrix},\quad A^2=-\omega^2I.
\]
指数级数的偶奇分解给出
\[
 e^{tA}=\begin{pmatrix}
 \cos(\omega t)&\sin(\omega t)/\omega\\
 -\omega\sin(\omega t)&\cos(\omega t)
 \end{pmatrix}.
\]
因此变常数公式的第一分量为
\[
 y_\nu(t)=\frac1\omega\int_0^t\sin(\omega(t-s))\cos(\nu s)\,\mathrm ds.
\]
积分，或验证方程及初值后使用唯一性，得到
\[
 y_\nu(t)=
 \begin{cases}
 \dfrac{\cos(\nu t)-\cos(\omega t)}{\omega^2-\nu^2},&\nu\ne\omega,\\[5pt]
 \dfrac{t}{2\omega}\sin(\omega t),&\nu=\omega.
 \end{cases}
\]
对任意 \(T>0\)，积分表达式和 \(|\cos(\nu s)-\cos(\omega s)|\leq|\nu-\omega|\cdot|s|\) 给出
\[
 \sup_{|t|\leq T}|y_\nu(t)-y_\omega(t)|
 \leq\frac{T^2}{2\omega}|\nu-\omega|.
\]
固定 \(\nu\ne\omega\) 时响应有界，但所得界 \(2/|\omega^2-\nu^2|\) 随 \(\nu\to\omega\) 增大。共振解则在 \(t=(\pi/2+2k\pi)/\omega\) 上线性增大，因而无界。紧时间连续依赖没有提供无穷时间的一致界。
\end{solution}

% 目标：用一个周期的演化寻找周期解，并证明吸引性与唯一性。
% 来源：自拟；不需要尚未建立的一般不动点或Floquet理论。
\begin{exercise}[周期外力的唯一周期响应]\label{b2:ex:ode-periodic-forcing}
设 \(\lambda,T>0\)，\(g\in C(\mathbb R;\mathbb R)\) 且 \(g(t+T)=g(t)\)。证明方程 \(x'=-\lambda x+g(t)\) 有且仅有一个 \(T\)-周期解 \(p\)，求出 \(p(0)\)，并证明任意其他解 \(x\) 满足
\[
 |x(t)-p(t)|=e^{-\lambda t}|x(0)-p(0)|\quad(t\geq0).
\]
\end{exercise}
\begin{solution}
由变常数公式，从 \(x(0)=a\) 出发的全时间解满足
\[
 x(T)=e^{-\lambda T}a+\int_0^T e^{-\lambda(T-s)}g(s)\,\mathrm ds.
\]
周期解必须满足 \(x(T)=a\)。由于 \(1-e^{-\lambda T}>0\)，唯一可能的初值是
\[
 a_* =\frac{\displaystyle\int_0^T e^{-\lambda(T-s)}g(s)\,\mathrm ds}
 {1-e^{-\lambda T}}.
\]
令 \(p\) 为这个初值的解。函数 \(t\mapsto p(t+T)\) 由外力的周期性满足同一个方程，且在 \(t=0\) 处取值 \(p(T)=a_*=p(0)\)。唯一性保证 \(p(t+T)=p(t)\) 对所有实数 \(t\) 成立，故必要条件也足够。任意两个解之差满足 \(z'=-\lambda z\)，从而得到题设的衰减恒等式。
\end{solution}

% 目标：检验流可微定理的C¹假设，区分Lipschitz连续依赖与可微依赖。
% 来源：自拟；正文没有以该流为证明前提。
\begin{exercise}[唯一的流未必可微]\label{b2:ex:ode-lipschitz-nonsmooth-flow}
求 \(x'=|x|\)、\(x(0)=a\in\mathbb R\) 的全时间流 \(\Phi_t(a)\)。证明对每个固定 \(t\)，该流关于 \(a\) 是 Lipschitz 映射；但当 \(t\ne0\) 时，它在 \(a=0\) 处不可微。说明这与本章哪些结论相容。
\end{exercise}
\begin{solution}
速度场全局 Lipschitz，且至多线性增长，因此解在整个实轴上唯一存在。分别解正、负半轴上的方程并核对初值，得到
\[
 \Phi_t(a)=
 \begin{cases}
 ae^t,&a\geq0,\\
 ae^{-t},&a<0.
 \end{cases}
\]
在每个半轴斜率的绝对值至多 \(e^{|t|}\)；跨过原点时把差分拆为两段，也得到 \(|\Phi_t(a)-\Phi_t(b)|\leq e^{|t|}|a-b|\)。然而原点的右导数为 \(e^t\)，左导数为 \(e^{-t}\)，二者恰在 \(t=0\) 时相等。局部 Lipschitz 条件保证唯一性和 Lipschitz 连续依赖，\cref{b2:thm:07-flow-derivative}还要求速度场属于 \(C^1\)；本例在原点不满足后一个条件。
\end{solution}

% 目标：计算参数变分，区分固定轨道段光滑依赖与最大生存时间变化。
% 来源：自拟；三种参数符号给出不同的最大区间。
\begin{exercise}[参数导数与存在时间]\label{b2:ex:ode-parameter-lifetime}
对实参数 \(\lambda\)，考虑 \(x'=x^2+\lambda\)、\(x(0)=0\)。分别求 \(\lambda>0\)、\(\lambda=0\)、\(\lambda<0\) 的最大解与最大时间区间。用参数变分方程求 \(\left.\partial_\lambda x(t,\lambda)\right|_{\lambda=0}\)，并解释为何 \(\lambda>0\) 时的有限时间爆破不违背 \(\lambda=0\) 附近的光滑依赖。
\end{exercise}
\begin{solution}
直接积分或代入验证可得
\[
 x(t,\lambda)=
 \begin{cases}
 \sqrt\lambda\tan(\sqrt\lambda\,t),&\lambda>0,\\
 0,&\lambda=0,\\
 -\sqrt{-\lambda}\tanh(\sqrt{-\lambda}\,t),&\lambda<0.
 \end{cases}
\]
当 \(\lambda>0\) 时最大区间为 \((-\pi/(2\sqrt\lambda),\pi/(2\sqrt\lambda))\)；另外两种情形的最大区间均为 \(\mathbb R\)。对于负参数，\(\tanh\) 形式也可写成指数的商，其分母处处为正。

由\cref{b2:prop:ode-parameter-variation}，固定任意紧时间区间，都可在 \(\lambda=0\) 的某邻域对解求参数导数。记 \(B(t)=\left.\partial_\lambda x(t,\lambda)\right|_{\lambda=0}\)，则参考解为零，变分方程为 \(B'=2\cdot0\cdot B+1\)、\(B(0)=0\)，故 \(B(t)=t\)。

给定 \(T>0\)，只要正参数满足 \(\lambda<\pi^2/(4T^2)\)，其解就存在于 \([-T,T]\)。因此任意固定紧时间段都保留了共同存在的参数邻域；当要求覆盖越来越长的时间段时，该邻域可以缩小。光滑依赖在最大解的开定义域上成立，不保证所有邻近参数共享同一个无限时间区间。
\end{solution}

% 目标：理解紧性构造的子列与近似方法的选解现象。
% 来源：自拟；在Peano正文非唯一模型上增加离散近似比较。
\begin{optionalexercise}[Euler 折线选择了哪一条解]\label{b2:ex:ode-peano-selection}
在 \([0,T]\)、\(T>0\) 上考虑 \(x'=2\sqrt{|x|}\)、\(x(0)=0\)。证明任意网格上从精确零初值出发的显式 Euler 折线都恒为零。另一方面，对 \(\varepsilon>0\)，求初值 \(x(0)=\varepsilon\) 的正向解，并证明它们当 \(\varepsilon\downarrow0\) 时一致趋于另一条零初值解。说明为何这不与 Peano 定理冲突。
\end{optionalexercise}
\begin{solution}
显式 Euler 节点递推为 \(x_{j+1}=x_j+2h_j\sqrt{|x_j|}\)。从 \(x_0=0\) 出发，归纳得到每个节点都为零，各段斜率也为零，所以整条折线恒为零，任意细化网格的极限仍为零解。

从正初值出发，解在正向时间保持正，由 \((\sqrt{x})'=1\) 得 \(x_\varepsilon(t)=(\sqrt\varepsilon+t)^2\)。在 \([0,T]\) 上，
\[
 \sup_{0\leq t\leq T}|x_\varepsilon(t)-t^2|
 \leq2T\sqrt\varepsilon+\varepsilon\longrightarrow0.
\]
极限 \(t^2\) 也满足零初值问题，却与零解不同。Peano 定理保证某条解可以由近似的紧性得到，没有断言所有近似过程都会选出同一条解，也没有给出初值唯一地确定的连续解映射。缺少唯一性时，不能把某个近似方案的极限当作全部解的描述。
\end{solution}

% 目标：结合旋转与径向衰减，区分渐近稳定和指数稳定。
% 来源：自拟；二维多项式耗散系统，避免重复正文一维模型。
\begin{optionalexercise}[旋转中的缓慢衰减]\label{b2:ex:ode-spiral-decay}
在 \(\mathbb R^2\) 上考虑
\[
 x'=-y-x(x^2+y^2),\quad y'=x-y(x^2+y^2).
\]
证明每条轨道在所有正向时间存在，原点渐近稳定，而且所有初值的解都趋于原点。求非零初值的半径与转角随时间的变化，并证明原点不满足局部指数稳定估计。
\end{optionalexercise}
\begin{solution}
取 \(V(x,y)=x^2+y^2\)。直接求导得 \(\dot V=-2(x^2+y^2)^2=-2V^2\)。因此轨道的半径不增加，在任意有限正向时间内留在固定闭球中，由延拓准则得到全正向存在。\(V\) 正定，\(\dot V\) 在原点外严格为负，故 Lyapunov 判据给出局部渐近稳定；下面的显式计算进一步给出所有初值的吸引性。

对非零初值，唯一性保证轨道在有限时刻不会碰到原点。令 \(r_0=\sqrt{x(0)^2+y(0)^2}>0\)。解 \(V'=-2V^2\) 得
\[
 r(t)=\frac{r_0}{\sqrt{1+2r_0^2t}},\quad t\geq0.
\]
沿轨道取连续的转角 \(\theta\)，则
\[
 \theta' =\frac{xy'-yx'}{x^2+y^2}=1,
 \quad\theta(t)=\theta(0)+t.
\]
轨道一边匀速转动，一边以代数速度靠近原点。半径不增加给出稳定性，任意 \(r_0\) 都有 \(r(t)\to0\)，故原点渐近稳定，且吸引范围是整个平面。

若存在 \(C,\gamma,\delta>0\) 使 \(r(t)\leq Ce^{-\gamma t}r_0\) 对所有 \(0<r_0<\delta\)、\(t\geq0\) 成立，固定其中一个 \(r_0\) 就会得到 \(e^{\gamma t}/\sqrt{1+2r_0^2t}\leq C\)。左侧趋于无穷，矛盾。因此渐近稳定没有自动带来指数速率。
\end{solution}
